\documentclass[10pt,letterpaper]{article}

\usepackage{ccn}
\usepackage{pslatex}

\usepackage[authoryear]{natbib}

\title{Modeling Nicotinic Receptor Dynamics in Alzheimer's Disease}

\author{
Valentin Jacquin under the supervision of Dr.\ Boris Gutkin \\
Group for Neural Theory \\
ENS Paris, France \\
}

% Add hyperref for clickable links in bibliography
\usepackage{hyperref}
\hypersetup{
    colorlinks=true,
    linkcolor=blue,
    citecolor=blue,
    filecolor=blue,
    urlcolor=blue
}

\begin{document}

\maketitle

\section{Abstract}
{\bf
% Provide a concise summary (150-250 words) of the planned
% internship project. This should include:
% • the scientific context,
% • the main research question,
% • the methodological approach,
% • the expected outcomes or observables.
Alzheimer's disease (AD) is a prevalent neurodegenerative disorder characterized by progressive cognitive decline. During its early stage, the $A\beta$ molecule accumulate and provoke various disruption in the neural system. Previous studies has explored the influence of its disruption on nicotinic acetylcholine receptors (nAChRs) on the hypofrontality observed in its early stage. Building upon their computational model of PFC, we aims to investigate how nAChR dysfunction affects PFC activity during working memory tasks, which are known to be impaired in AD. Incorporating nAChR modulation, the model will allows us to explore the impact of receptor dysfunction on network bistability and transition dynamics. By implementing a ring attractor architecture, we will simulate working memory processes and assess how impairments in specific nAChR subtypes influence the network's ability to maintain stable activity states. We will analyze firing rates, transition frequencies, and working memory performance metrics to elucidate the relationship between nAChR dysfunction and cognitive deficits in AD.
}
\section{Scientific Background and Context}
{
% Summarize the key scientific context supporting your future
% internship project, based on the bibliography you have
% explored during the pre-internship period. This section
% should:
% • briefly introduce the scientific domain (neuroscience, 
% cognitive science, AI, computational modelling, etc.),
% • highlight the main relevant theories or previous findings,
% • show that you have understood the state of the art around 
% your planned topic.
Alzheimer's disease (AD) is a widespread neurodegenerative disorder with a high prevalence in elderly population, currently affecting over 50 millions individuals worldwide~\citep{Breijyeh2020}. However, the origins of the cognitive impairments associated with AD remain poorly understood. While \citet{Palop2007} have reported increases in cortical activity in mice with AD, the underlying % i think i should add another studies that talk about its hypoactivity to be like -> there's a discordance in the literature, which shows that we still don't understand the mechanisms
mechanisms remain debated. Experimental evidence indicates that amyloid-beta (A$\beta$) peptides modulate neuron's signals by either activating or inhibiting nicotinic acetylcholine receptor (nAChR). More specifically, $\alpha7$ and $\beta2^*$ nAChR's subunits seems to be particularly affected by A$\beta$ interactions~\citep{Lombardo2015,Parri2011}.

In this work, we hypothesize that the progressive accumulation of A$\beta$ disrupts nAChR mechanisms, leading to a disregulation in prefrontal cortex (PFC) network dynamics. Thus, contributing to the cognitive deficits observed in AD. As nAChRs are only express in interneurons in layer II/III~\citep{Poorthuis2013}, nicotinic signaling is expected to impact pyramidal cells through interneurons modulation. In this perspective, study from \citet{Koukouli2017} have demonstrated the link betsween nAChRs disruption and pyramidal neurons activity, underlying the interneurons implication in PFC network dynamics. The specific contribution of different receptor subunits to this modulation has also been highlighted~\citep{Maskos2014}. Still, the precise interactions between nAChR dysfunction and the resulting alterations in PFC network dynamics in the context of AD remain to be explained.

To better capture these mechanisms, previous work developed a reduced population-level model describing interactions between pyramidal neurons and major interneuron subtypes within the PFC, constrained by biologically observed data~\citep{Koukouli2025}. This framework reproduces the hyperactivity observed in early AD stages and identifies $\alpha7$ nAChRs as critical modulators of this phenomenon. Moreover, it suggests that $\beta2^*$ and $\alpha5$ nAChRs could represent promising therapeutic targets for restoring physiological network balance.

}
\begin{quote}
\small
\textbf{Keywords:} Alzheimer's disease, nicotinic acetylcholine receptors, prefrontal cortex, computational modeling, working memory.
\end{quote}

\section{Research Question and Working Hypotheses}
% Formulate precisely:
% • the main research question you intend to address,
% • any sub-questions if relevant,
% • the hypotheses you plan to test during the internship.
% This part must make explicit what you aim to discover, measure, or demonstrate.

While previous explorations have showned promising results in linking nAChR dysfunction to altered PFC activity patterns~\citep{Rooy2021,Koukouli2025} they primarily focused on resting-state dynamics. However, to link these findings to the observed cognitive deficits in AD, it is essential to investigate how these network alterations manifest during active cognitive tasks.
The PFC is known to be a key region in working memory, which deteriorate in early AD~\citep{Huntley2010}. Furthermore, electrophysiological recordings during working memory tasks have been shown to reliably distinguish AD patients from healthy controls~\citep{Perpetuini2020}, underscoring the connection between altered neural dynamics and cognitive impairment.

Therefore, working memory task appears as a good candidate to explore how nAChR dysfunction affects PFC network dynamics in an active context. We will investigate how disruptions in the different nAChR subtypes influence the network's ability to switch between two self-sustained states, which are thought to underlie working memory processes~\citep{Compte2003}. We hypothesize that impairments in specific nAChR subtypes will lead to deficits in the stability and transition dynamics between these states, thereby contributing to the working memory impairments observed in AD.

To another extent, we will also explore how varying levels of external input or sensitivity on specific nAChR subtypes affect these dynamics. This, in the aim to provides insights into potential therapeutic strategies that could restore normal network function and improve cognitive outcomes in AD patients.

\section{Planned Methodology}
% Describe the methodological strategy you expect to follow. 
% Depending on your project, this may include:
% • computational methods or algorithms,
% • experimental paradigms or behavioral tasks,
% • modelling tools or simulations,
% • data analysis techniques.
% This section should emphasize feasibility within the M2 internship timeframe.

To tackle this question, we will build upon the previously developped nAChR-modulated PFC network model described in \citet{Koukouli2025}.

\subsection{Modulation of parameters affecting bistability}
Previous results have shown an oscillation between high and low activity states in the PFC~\citep{Rooy2021}, but at another order of magnitude than the one observed experimentally during working memory tasks~\citep{Compte2003}. Thus, it is needed to bridge the gap closer to experimental observations, as we aim to maintain a biologically relevant model. To do so, we will modulate current inputs as we expect them to highly affect the characteristics of both states.

\subsection{Working memory architecture}
Once we achieve a network behavior consistent with experimental findings, we will build a ring attractor network. This architecture has been widely used to model working memory processes, notably in~\citep{Wang2004} where the authors also modelize the role of interneurons. However, they first do not consider nAChR modulation, and second use a arbitraty connectum framework between nodes. 

Here, we will implement a similar architecture but defining each node as a local PFC circuit described by our nAChR-modulated model. We will consider PYR to PYR connections as the driver of the excitatory connectum, while PV interneurons will mediate the inhibitory connections between nodes, as they are both known to be able to project horizontally~\citep{Levitt1993,Druga2023}. SOM and VIP interneurons will be considered as local modulators of the network dynamics as they mainly project locally~\citep{Karnani2016,Riedemann2019}.

\section{Expected Observables and Data}
% Explain the type of observables you expect to analyse or generate.
% Examples include:
% • behavioral responses or psychophysical measurements,
% • neural signals or imaging data,
% • model outputs or synthetic datasets,
% • statistical indicators or quantitative metrics.
% If real data will be used, specify their nature and origin (public datasets, provided by the host lab, etc.).

We will primarily focus on analyzing the firing rates of the different neuron populations (PYR, PV, SOM, VIP) within the PFC network model. The following observables will be used to assess the impact of nAChR dysfunction and input modulation on network dynamics.

\subsection{Characteristics of activity states}
As previously mention, we first aim to get closer to experimental observations during working memory tasks~\citep{Compte2003}. We will quantify the mean firing rates and variance of both high and low-activity states along to varying input current levels. We expect that increasing input current will lead to higher firing rates in both states.

\subsection{Transition dynamics}
We will quantify the frequency of transitions between high and low-activity states. This should be highly modulated by the input current, and we expect that stronger inputs may lead to prolonged high-activity states, potentially resulting in a one-stability state at a certain threshold. Whereas, lower inputs may favor a more stable low-activity state.

\subsection{Working memory performance metrics}
In the ring attractor model, we will first assess the network's ability to form and maintain a bump of activity in response to a transient stimulus. Then, we will measure the population vectors to quantify the accuracy of working memory representations. We expect the network to be able to maintain stable activity bumps, sucessflully encoding the stimulus information. In the context of nAChR dysfunction, we will evaluate how impairments in specific receptor subtypes affect the stability and fidelity of these representations.

\bibliographystyle{plainnat}
\bibliography{references}

\end{document}